\exercisesection{\S~4}

The notations and assumptions are those of nos. 1, 2, 3 of §4.

\begin{enumerate}
\def\labelenumi{\arabic{enumi})}
\tightlist
\item
  Take \(\mathrm{G} = \mathbf{GL}_{n}(k), n \geq 0.\)
\end{enumerate}

\begin{enumerate}
\def\labelenumi{\alph{enumi})}
\item
  Show that \(r_{\mathrm{Ad}}^0 (g) = n\) for all \(g \in \mathrm{G}\) .
\item
  Show that an element \(g \in G\) is regular if and only if its characteristic polynomial \(\mathrm{P}_{g}(\mathrm{T}) = \det(\mathrm{T} - g)\) is separable; this is equivalent to saying that the discriminant (Algebra, Chap. IV, §1, no. 10) of \(\mathrm{P}_{g}(\mathrm{T})\) is \(\neq 0\) .
\end{enumerate}

\begin{enumerate}
\def\labelenumi{\arabic{enumi})}
\setcounter{enumi}{1}
\item
  Construct a Lie group G such that the function \(r_{Ad}^{0}\) is not constant. (Take g abelian \(\neq 0\) and Ad non-trivial.)
\item
  Let \((\rho_{i})_{i\in I}\) be a countable family of analytic linear representations of G. Prove that the elements of G which are regular for all the \(\rho_{i}\) constitute a dense subset of G. Construct an example showing that the countability assumption cannot be omitted.
\item
  Assume that \(k = \mathbf{C}\) and that \(\mathbf{G}\) is connected. Prove the equivalence of the following properties:
\end{enumerate}

\begin{enumerate}
\def\labelenumi{(\roman{enumi})}
\item
  G is nilpotent.
\item
  Every element \(\neq 1\) of \(G\) is regular.
\end{enumerate}

(Show first that (ii) implies

\begin{enumerate}
\def\labelenumi{(\roman{enumi})}
\setcounter{enumi}{1}
\tightlist
\item
  \(^{\prime}\) Every element \(\neq 0\) of g is regular.
\end{enumerate}

Remark next that, if \(g \neq 0\) , then g contains elements \(x \neq 0\) such that ad x is nilpotent, cf.~§3, Exerc. 10. Deduce that g is nilpotent, hence (i).

